% ==============================================================================
% PREAMBLE: Formato oficial Trabajo de Grado - Instituto de Física, UdeA
% ==============================================================================
\usepackage[utf8]{inputenc}
% 'es-noquoting' desactiva la activación de < y > que rompe TikZ
\usepackage[spanish,es-nodecimaldot,es-tabla,es-noquoting]{babel}
\spanishplainpercent

% Geometría de página (Márgenes estándar UdeA)
\usepackage[letterpaper, top=3.0cm, bottom=2.5cm, left=3.0cm, right=2.5cm]{geometry}

% Tipografía idéntica a la tesis de referencia (Times / NewTX)
\IfFileExists{newtxtext.sty}{
    \usepackage{newtxtext}
    \usepackage{newtxmath}
}{
    \usepackage{mathptmx}
}
\usepackage{microtype}

% Espaciado e interlineado
\usepackage{setspace}
\setstretch{1.15}
\setlength{\parskip}{6pt}
\setlength{\parindent}{0pt}

% Enlaces y URLs con división automática
\usepackage{xcolor}
\definecolor{UdeAGreen}{HTML}{006837} % Verde institucional UdeA
\usepackage{xurl}
\usepackage{hyperref}
\hypersetup{
    colorlinks=true,
    linkcolor=black,
    citecolor=black,
    urlcolor=UdeAGreen
}

% Encabezados y pies de página
\usepackage{fancyhdr}
\pagestyle{fancy}
\fancyhf{}
\fancyhead[L]{\small\textit{WALO: Localización de Fuentes de Ondas Acústicas}}\fancyhead[R]{\small\thepage}
\renewcommand{\headrulewidth}{0.4pt}
\renewcommand{\footrulewidth}{0pt}
\setlength{\headheight}{15pt}

\fancypagestyle{plain}{
    \fancyhf{}
    \fancyhead[R]{\small\thepage}
    \renewcommand{\headrulewidth}{0pt}
}

% Formato sobrio de títulos
\usepackage{titlesec}
\titleformat{\section}{\large\bfseries}{\thesection.}{0.5em}{}
\titleformat{\subsection}{\normalsize\bfseries}{\thesubsection.}{0.5em}{}
\titleformat{\subsubsection}{\small\bfseries}{\thesubsubsection.}{0.5em}{}

% Matemáticas y física
\let\Bbbk\relax
\usepackage{amsmath,amssymb,amsfonts,bm}

\IfFileExists{physics.sty}{
    \usepackage{physics}
}{
    \newcommand{\norm}[1]{\left\lVert#1\right\rVert}
    \newcommand{\abs}[1]{\left\lvert#1\right\rvert}
}

% Unidades físicas
\IfFileExists{siunitx.sty}{
    \usepackage{siunitx}
    \sisetup{
        output-decimal-marker = {.},
        per-mode = symbol
    }
}{}

% Figuras, tablas y TikZ (con compatibilidad de babel activada)
\usepackage{graphicx}
\usepackage{booktabs}
\usepackage{tabularx}
\usepackage{tikz}
\usetikzlibrary{shapes.geometric, arrows.meta, positioning, calc, babel, patterns}
\usepackage{caption}
\captionsetup{font=small, labelfont=bf, labelsep=period}

% Citas bibliográficas estilo IEEE
\usepackage[numbers,sort&compress]{natbib}
\bibliographystyle{ieeetr}

% ==============================================================================
% OPTIMIZACIÓN DE PARÁMETROS DE FLOTACIÓN (Evita páginas con figuras aisladas)
% ==============================================================================
\renewcommand{\topfraction}{0.85}       % Máx. proporción de página ocupada por figuras arriba (def: 0.70)
\renewcommand{\bottomfraction}{0.80}    % Máx. proporción ocupada por figuras abajo (def: 0.30)
\renewcommand{\textfraction}{0.12}      % Mín. proporción obligatoria de texto por página (def: 0.20)
\renewcommand{\floatpagefraction}{0.75} % Solo crea página de figuras si la figura ocupa al menos 75% (def: 0.50)
\setcounter{topnumber}{3}               % Máximo 3 figuras en la parte superior
\setcounter{bottomnumber}{2}            % Máximo 2 figuras en la parte inferior
\setcounter{totalnumber}{4}             % Máximo 4 flotantes por página